% Môn: Vật lý
% Lớp: 12
% Chương: Chương I. Vật lí nhiệt
% Bài: L12C1 Bài 6. Nhiệt hóa hơi riêng · Bài toán thực nghiệm xác định nhiệt hóa hơi riêng của nước thông qua các phép đo thời gian, công suất ấm đun và độ hụt k
% ===== Câu 96 =====
% ID: L12C1B6-182-TN
% Mức: NB
\dangbt{Bài toán thực nghiệm xác định nhiệt hóa hơi riêng của nước thông qua các phép đo thời gian, công suất ấm đun và độ hụt k}
\begin{ex}
% ID: L12C1B6-182-TN
% Mức: NB
Nhiệt dung riêng của một chất là nhiệt lượng cần cung cấp để nhiệt độ của
\choice
{\True $1\ \mathrm{kg}$ chất đó tăng thêm $1\ \mathrm{K}$}
{$1$ phân tử chất đó tăng thêm $1K$}
{$1\ \mathrm{m^3}$ chất đó tăng thêm $1\ \mathrm{K}$}
{$1\ \mathrm{mol}$ chất đó tăng thêm $1\ \mathrm{K}$}
\loigiai{
Nhiệt dung riêng của một chất là nhiệt lượng cần cung cấp để nhiệt độ của $1\,\text{kg}$ chất đó tăng thêm 1 K.
}
\end{ex}

% ===== Câu 100 =====
% ID: L12C1B6-183-TN
% Mức: NB
\begin{ex}
% ID: L12C1B6-183-TN
% Mức: NB
Trong hệ SI, đơn vị của nhiệt dung riêng là
\choice
{\True $J/kg.K$}
{$J$}
{$J/kg$}
{$J/ K$}
\loigiai{
Nhiệt dung riêng $c$ của một chất được xác định từ công thức
  Q=mc $\Delta$ t
 $\quad\Longrightarrow\quad$ c= $\frac{Q}{m\Delta t}$. 
 
 Trong hệ SI:
 
 
• Nhiệt lượng $Q$ có đơn vị là $J$ (Jun);
 
• Khối lượng $m$ có đơn vị là $kg$;
 
• Độ tăng nhiệt độ $\Delta t$ có đơn vị là $K$.
 
 
 Do đó, đơn vị của nhiệt dung riêng là
 $[c]=\dfrac{J}{kg\cdot K}.$
 
 Vậy đáp án đúng là $J/kg.K$.
}
\end{ex}

% ===== Câu 104 =====
% ID: L12C1B6-184-TN
% Mức: NB
\begin{ex}
% ID: L12C1B6-184-TN
% Mức: NB
Trong một bài thực hành vật lí, một nhóm học sinh đun nóng một khối lượng nước m = $0,136\,\text{kg}$ bằng một dây điện trở có công suất tỏa nhiệt không đổi $P$ = $18,2\,\text{W}$. Nhiệt độ ban đầu của nước là 27 độ C. Sau thời gian đun liên tục $\Delta t = 900\,\text{s}$, nhiệt kế chỉ nhiệt độ nước là 54 độ C. Bỏ qua mọi hao phí nhiệt ra môi trường và nhiệt lượng làm nóng bình đun. Hãy tính giá trị thực nghiệm của nhiệt dung riêng của nước thu được từ phép đo này.
\choice
{$4180\,\text{J}$ /kg. $K$}
{\True $4461\,\text{J}$ /kg. $K$}
{$3985\,\text{J}$ /kg. $K$}
{$4250\,\text{J}$ /kg. $K$}
\loigiai{
Nhiệt lượng do dây điện trở tỏa ra trong thời gian đun là:
$Q_{\text{tỏa}} = P\Delta t = 18,2 \cdot 900 = 16380 \mathrm{J}.$

Độ tăng nhiệt độ của lượng nước là:
$\Delta T = 54 - 27 = 27 \mathrm{K}.$

Coi toàn bộ nhiệt lượng tỏa ra được nước hấp thụ hoàn toàn:
$P\Delta t = mc\Delta T.$

Nhiệt dung riêng của nước đo được bằng thực nghiệm là:
$c = \dfrac{P\Delta t}{m\Delta T}.$

Thay số:
$c = \dfrac{16380}{0,136 \cdot 27} = \dfrac{16380}{3,672} \approx 4461 \mathrm{J/(kg\cdot K)}.$
}
\end{ex}

% ===== Câu 108 =====
% ID: L12C1B6-185-TN
% Mức: TH
\begin{ex}
% ID: L12C1B6-185-TN
% Mức: TH
Một phép đo thực nghiệm đo nhiệt dung riêng của nước thu được đồ thị biểu diễn sự thay đổi nhiệt độ theo thời gian đun như hình vẽ dưới dạng một đường thẳng. Biết công suất đun của nguồn điện là $P$, khối lượng nước là m. Hệ số góc của đường thẳng này được xác định bằng biểu thức nào sau đây?
\choice
{$k=\dfrac{Pm}{c}$}
{$k=\dfrac{c}{Pm}$}
{\True $k=\dfrac{P}{mc}$}
{$k=\dfrac{mc}{P}$}
\loigiai{
Ta có biểu thức bảo toàn năng lượng khi bỏ qua hao phí:
$mc(T-T_0)=Pt.$

Suy ra:
$T = T_0 + \dfrac{P}{mc}t.$

Đây là phương trình đường thẳng có dạng:
$y=ax+b.$

Với trục tung là nhiệt độ $T$ và trục hoành là thời gian $t$, hệ số góc của đồ thị là:
$k=\dfrac{P}{mc}.$
}
\end{ex}

% ===== Câu 112 =====
% ID: L12C1B6-186-TN
% Mức: TH
\begin{ex}
% ID: L12C1B6-186-TN
% Mức: TH
Một thợ rèn nhúng một con dao rựa bằng thép có khối lượng $1,1\,\text{kg}$ đang ở nhiệt độ nóng đỏ là 850 độ $C$ vào trong một bể chứa 200 lít nước ở nhiệt độ ngoài trời là 27 độ $C$ để làm nguội nhanh. Giả sử không có nước bị bay hơi và bỏ qua sự truyền nhiệt cho thành bể cũng như môi trường bên ngoài. Biết nhiệt dung riêng của thép là $460\,\text{J}$ /kg. $K$, của nước là $4180\,\text{J}$ /kg. $K$ và khối lượng riêng của nước là $1\,\text{kg}$ /lít. Nhiệt độ của nước trong bể khi đạt trạng thái cân bằng nhiệt xấp xĩ bằng bao nhiêu?
\choice
{27,1 độ $C$}
{\True 27,5 độ $C$}
{28,2 độ $C$}
{29,5 độ $C$}
\loigiai{
Khối lượng nước trong bể là:
$m_2 = 200 \cdot 1 = 200 \mathrm{kg}.$

Gọi $T$ là nhiệt độ của hệ khi đạt trạng thái cân bằng nhiệt.

Nhiệt lượng do con dao thép tỏa ra khi nguội từ $850^\circ\mathrm{C}$ xuống $T$ là:
$Q_{\text{tỏa}} = m_1c_1(T_1-T).$

Thay số:
$Q_{\text{tỏa}} = 1,1 \cdot 460(850-T) = 506(850-T) = 430100 - 506T.$

Nhiệt lượng do nước thu vào để tăng nhiệt độ từ $27^\circ\mathrm{C}$ lên $T$ là:
$Q_{\text{thu}} = m_2c_2(T-T_2).$

Thay số:
$Q_{\text{thu}} = 200 \cdot 4180(T-27) = 836000(T-27) = 836000T - 22572000.$

Theo phương trình cân bằng nhiệt:
$Q_{\text{tỏa}} = Q_{\text{thu}}.$

Suy ra:
$430100 - 506T = 836000T - 22572000.$

Do đó:
$836506T = 23002100.$

Vậy:
$T \approx 27,5^\circ\mathrm{C}.$
}
\end{ex}

% ===== Câu 116 =====
% ID: L12C1B6-187-TN
% Mức: VD
\begin{ex}
% ID: L12C1B6-187-TN
% Mức: VD
Một bình hình trụ có bán kính đáy $R_1=20\ \mathrm{cm}$, được đặt thẳng đứng và chứa nước ở nhiệt độ $T_1=20\ ^\circ\mathrm{C}$. Người ta thả vào bình một quả cầu bằng nhôm có bán kính $R_2=10\ \mathrm{cm}$, ở nhiệt độ $T_2=40\ ^\circ\mathrm{C}$. Khi cân bằng, quả cầu nằm dưới đáy bình và mực nước đi qua tâm quả cầu. Biết khối lượng riêng của nước và nhôm lần lượt là $\rho_1=1,000\ \mathrm{kg/m^3}$ và $\rho_2=2,700\ \mathrm{kg/m^3}$; nhiệt dung riêng của nước và nhôm lần lượt là $c_1=4,200\ \mathrm{J/(kg\cdot K)}$ và $c_2=880\ \mathrm{J/(kg\cdot K)}$. Bỏ qua sự trao đổi nhiệt với bình và môi trường. Nhiệt độ của nước khi cân bằng nhiệt là
\choice
{$20,7\ ^\circ\mathrm{C}$}
{$24,8\ ^\circ\mathrm{C}$}
{$23,95\ ^\circ\mathrm{C}$}
{\True $23,7\ ^\circ\mathrm{C}$}
\loigiai{
Vì $\rho_2\gt \rho_1$ nên quả cầu nhôm chìm và nằm dưới đáy bình.

 Thể tích của quả cầu nhôm là
  \begin{aligned}
 V_2
 &= $\dfrac{4}{3}\pi$ R_2^3
&= $\dfrac{4}{3}\pi\cdot(0,1)^3$ &= $\dfrac{4}{3}\pi\cdot$ 10^{-3} $\\mathrm{m^3}$.
 \end{aligned} 
 
 Khi mực nước đi qua tâm quả cầu, phần quả cầu ngập trong nước là một nửa quả cầu. Do đó, thể tích nước trong bình là
  \begin{aligned}
 V_1
 &= $\pi$ R_1^2R_2- $\dfrac{1}{2}V_2$ &= $\pi\cdot(0,2)^2\cdot$ 0,1
 - $\dfrac{1}{2}\cdot\dfrac{4}{3}\pi\cdot$ 10^{-3}
&= $\dfrac{10}{3}\pi\cdot$ 10^{-3} $\\mathrm{m^3}$.
 \end{aligned} 
 
 Gọi $T$ là nhiệt độ cân bằng của hệ. Nhiệt lượng quả cầu nhôm tỏa ra là
  \begin{aligned}
 Q_{\mathrm{tỏa}}
 &=m_2\cdot c_2 $\cdot(T_2-T)$ &= $\rho_2V_2\cdot$ c_2 $\cdot(40-T)$ &=2 700 $\cdot\dfrac{4}{3}\pi\cdot$ 10^{-3}
 \cdot 880 $\cdot(40-T)$ &=3 168 $\pi(40-T)\\mathrm{J}$.
 \end{aligned} 
 
 Nhiệt lượng nước thu vào là
  \begin{aligned}
 Q_{ $\mathrm{thu}$ }
 &=m_1\cdot c_1 $\cdot(T-T_1)$ &= $\rho_1V_1\cdot$ c_1 $\cdot(T-20)$ &=1 000 $\cdot\dfrac{10}{3}\pi\cdot$ 10^{-3}
 \cdot 4 200 $\cdot(T-20)$ &=14 000 $\pi(T-20)\\mathrm{J}$.
 \end{aligned} 
 
 Áp dụng phương trình cân bằng nhiệt:
  \begin{aligned}
 Q_{\mathrm{tỏa}}&=Q_{ $\mathrm{thu}$ },
3 168 $\pi(40-T)$ &=14 000 $\pi(T-20),$ T&\approx 23,7 $\$ ^ $\circ\mathrm{C}$.
 \end{aligned}
}
\end{ex}

% ===== Câu 120 =====
% ID: L12C1B6-188-TN
% Mức: VD
\begin{ex}
% ID: L12C1B6-188-TN
% Mức: VD
Đổ $50\ \mathrm{g}$ chất lỏng $A$ có nhiệt dung riêng $c_{A}$ ở $70^{\circ}\mathrm{C}$ vào $100\ \mathrm{g}$ chất lỏng $B$ có nhiệt dung riêng $c_{B}$ ở nhiệt độ $100^{\circ}\mathrm{C}$ thì nhiệt độ cân bằng của hỗn hợp là $90^{\circ}\mathrm{C}$. Bỏ qua mọi hao phí. Tỉ số $\dfrac{c_{A}}{c_{B}}$ bằng
\choice
{\True $1$}
{$2$}
{$4$}
{$3$}
\loigiai{
Chất lỏng $A$ có nhiệt độ ban đầu thấp hơn nhiệt độ cân bằng nên $A$ \textbf{thu nhiệt}; chất lỏng $B$ có nhiệt độ ban đầu cao hơn nhiệt độ cân bằng nên $B$ \textbf{tỏa nhiệt}.
 
 $\textbf$ {Nhiệt lượng chất lỏng $A$ thu vào} khi tăng từ $70^\circ\mathrm{C}$ lên $90^\circ\mathrm{C}$:
 $Q_{\text{thu}}=m_Ac_A\left(90-70\right)=50c_A\cdot20=1000c_A.\textbf$ {Nhiệt lượng chất lỏngB $tỏa ra} khi giảm từ$ 100^ $\circ\mathrm{C}$ xuống $90^$ \circ\mathrm{C} $:$ Q_{\text{tỏa}}=m_Bc_ $B$ \left(100-90\right)=100c_B\cdot10=1000c_B $.
 
 Theo phương trình cân bằng nhiệt:$Q_{$\text{thu}$}=Q_{\text{tỏa}}$, ta có
  1000c_A=1000c_B
 \quad\Longrightarrow\quad
 \frac{c_A}{c_B}=1. 
 
 Vậy tỉ số$ \dfrac{c_A}{c_B}=1.
}
\end{ex}

% ===== Câu 124 =====
% ID: L12C1B6-189-TN
% Mức: VD
\begin{ex}
% ID: L12C1B6-189-TN
% Mức: VD
Một bình cách nhiệt được ngăn thành hai phần bằng nhau bằng một vách ngăn. Hai phần của bình chứa hai chất lỏng có khối lượng lần lượt là $m_1$, $m_2$; nhiệt dung riêng lần lượt là $c_1$, $c_2$ và nhiệt độ ban đầu lần lượt là $t_1$, $t_2$, với $t_1\gt t_2$. Khi bỏ vách ngăn, hỗn hợp hai chất lỏng đạt nhiệt độ cân bằng $t$. Biết
	$t_1-t=\dfrac{1}{2}(t_1-t_2).$
	Tỉ số $\dfrac{m_1}{m_2}$ có giá trị là
\choice
{$\dfrac{m_1}{m_2}=1+\dfrac{c_2}{c_1}$}
{\True $\dfrac{m_1}{m_2}=\dfrac{c_2}{c_1}$}
{$\dfrac{m_1}{m_2}=\dfrac{c_1}{c_2}$}
{$\dfrac{m_1}{m_2}=1+\dfrac{c_1}{c_2}$}
\loigiai{
Từ giả thiết, ta có
  \begin{aligned}
 t_1-t
 &= $\dfrac{1}{2}(t_1-t_2),\$ 2 $(t_1-t)$ &=t_1-t_2, $\$
 t_1-t
 &=t-t_2.
 \end{aligned} 
 

 Vì bình cách nhiệt nên nhiệt lượng chất lỏng thứ nhất tỏa ra bằng nhiệt lượng chất lỏng thứ hai thu vào:
  \begin{aligned}
 Q_{\mathrm{tỏa}}&=Q_{ $\mathrm{thu}$ },
m_1\cdot c_1 $\cdot(t_1-t)$ &=m_2\cdot c_2 $\cdot(t-t_2)$.
 \end{aligned} 
 
 Do $t_1-t=t-t_2$, suy ra
  \begin{aligned}
 m_1c_1&=m_2c_2,
$\dfrac{m_1}{m_2}$ &= $\dfrac{c_2}{c_1}$.
 \end{aligned}
}
\end{ex}

% ===== Câu 128 =====
% ID: L12C1B6-190-TN
% Mức: VD
\begin{ex}
% ID: L12C1B6-190-TN
% Mức: VD
Người ta thả một vật rắn khối lượng $m_1$, nhiệt độ $150^\circ \mathrm{C}$ vào một bình chứa, bên trong chứa nước có khối lượng $m_2$. Khi cân bằng nhiệt, nhiệt độ của nước tăng từ $10^\circ \mathrm{C}$ đến $50^\circ \mathrm{C}$. Gọi $c_1,\ c_2$ lần lượt là nhiệt dung riêng của vật rắn và nhiệt dung riêng của nước. Bỏ qua sự hấp thụ nhiệt của bình và môi trường xung quanh. Hệ thức nào sau đây là đúng?
\choice
{$\dfrac{m_1c_1}{m_2c_2}=\dfrac{7}{2}$}
{$\dfrac{m_1c_1}{m_2c_2}=\dfrac{2}{7}$}
{$\dfrac{m_1c_1}{m_2c_2}=\dfrac{5}{2}$}
{\True $\dfrac{m_1c_1}{m_2c_2}=\dfrac{2}{5}$}
\loigiai{
Nhiệt độ cân bằng của hệ là $t=50^\circ\mathrm{C}$.
 
 \textbf{Nhiệt lượng vật rắn tỏa ra} khi hạ nhiệt từ $150^\circ\mathrm{C}$ xuống $50^\circ\mathrm{C}$:
 $Q_{\text{tỏa}}=m_1c_1\left(150-50\right)=100 m_1c_1.$ \textbf{Nhiệt lượng nước thu vào} khi tăng nhiệt từ10^ $\circ\mathrm{C}$ lên $50^$ \circ\mathrm{C} $:$ Q_{ $\text{thu}$ }=m_2c_2 $\left(50-10\right)=40$ m_2c_2. $Theo phương trình cân bằng nhiệt:$ Q_{\text{tỏa}}=Q_{ $\text{thu}$ } $, ta có
  100 m_1c_1=40 m_2c_2
 \quad\Longrightarrow\quad
 \frac{m_1c_1}{m_2c_2}=\frac{40}{100}=\frac{2}{5}. 
 
 Vậy hệ thức đúng là$ \dfrac{m_1c_1}{m_2c_2}=\dfrac{2}{5}.
}
\end{ex}
